The diagnostics in \texttt{results/} identify where the loss originates.
\begin{enumerate}
\item \emph{Not the Jacobian, not quadrature.} The loss occurs on meshes with
  $\det\bm J>0$ ($\alpha=30^\circ$) and persists with exact integration
  (\cref{sec:quad,sec:poly}).
\item \emph{Not consistency.} The residual of the exact nodal values,
  $\bm r=\bm K\bm u_I-\bm F$, is \emph{smaller} for \PM\ than for \PP\ in all eleven cases
  tested (\texttt{residual\_smooth.txt}; e.g.\ $\alpha=30^\circ$, $N=8$, $p=4$:
  $3.0\cdot10^{-5}$ versus $3.7\cdot10^{-4}$), whereas the \PM\ nodal error is $9$ times
  larger ($2.8\cdot10^{-5}$ versus $3.0\cdot10^{-6}$). The ratio of error to residual of
  \PM\ grows under refinement ($0.05$ at $N=4$, $0.9$ at $N=8$).
\item \emph{Stability of the unsymmetric operator.} With $\bm u_h-\bm u_I=-\bm K^{-1}\bm r=
  -\sum_k\bm v_k(\bm u_k^{\mathsf T}\bm r)/\sigma_k$ (singular triplets of $\bm K$;
  \texttt{spectrum\_smooth.txt}, smooth field), the error of \PP\ and of \PM\ for
  $\alpha\le15^\circ$ is spread over the spectrum (largest single contribution $2$--$8\,\%$ of
  $\|\bm u_h-\bm u_I\|^2$, from mid-spectrum modes). For $\alpha\ge30^\circ$, $p=4$, \PM\ has
  singular values below the smallest singular value of \PP\ (down to $0.17\,\sigma_1^{\PP}$),
  and in four of the five cases tested a single one of the lowest six triplets carries
  $21$--$68\,\%$ of the error (the fifth, $\alpha=30^\circ$, $N=4$, $p=6$: triplet $k=28$,
  $9\,\%$). The
  residual of the non-conforming Petrov--Galerkin scheme is not orthogonal to these
  low modes, unlike the Galerkin residual of \PP, so they are excited and amplified.
\item \emph{Ill-conditioned metric interpolation.} The Lebesgue constant
  $\Lambda_M=\max_{\bx\in\Omega^e}\sum_i|M_i^e(\bx)|$ of the metric functions (basis
  independent) equals the GLL value on affine elements, but grows by orders of magnitude
  with $\alpha$ and $p$ (\cref{tab:leb}): $\Lambda_M\le163$ for $\alpha=15^\circ$, $p\le6$,
  but $1.1\cdot10^3$ ($\alpha=30^\circ$, $p=4$), $9\cdot10^5$ ($45^\circ$, $p=5$) and
  $5\cdot10^{12}$ ($60^\circ$, $p=8$). The element condition numbers of \PM\ exceed those of
  \PP\ by similar factors (\cref{tab:cond}). The images of the GLL nodes of strongly
  non-affine elements come close to configurations that are not unisolvent for $Q_p$ in
  skew coordinates (``interpolation failure'' in \citep{Rajendran2010technique}); the
  metric trial functions then have large gradients, the bilinear form loses its stability
  constant, and round-off is amplified (\cref{sec:poly}). All cases of the benign regime
  of \cref{sec:conv} have $\Lambda_M\le2.5\cdot10^2$, all degraded cases
  $\Lambda_M\ge10^3$; the converse does not hold strictly ($\alpha=30^\circ$, $p=6$:
  $\Lambda_M=1.1\cdot10^4$, yet \PM\ is $500$ times more accurate than \PP\ at $N=8$), so
  $\Lambda_M$ is a necessary rather than a sufficient indicator. The maximum over the mesh does not decrease under
  refinement because the distortion is self-similar; it varies with $N$ because the
  element shapes near the curved boundaries change.
\end{enumerate}

\begin{table}[t]
  \centering\small
  \begin{minipage}[t]{0.48\linewidth}\centering
    \captionof{table}{Lebesgue constant $\Lambda_M$ of the metric interpolation, maximum over
      the elements, $N=8$ (GLL: affine element).}\label{tab:leb}
    \resizebox{\linewidth}{!}{\begin{tabular}{lrrrrrrr}
\toprule
& \multicolumn{7}{c}{$p$}\\
\cmidrule(lr){2-8}
 & 2 & 3 & 4 & 5 & 6 & 7 & 8\\
\midrule
\multicolumn{8}{l}{\PM\ (unsymmetric, skew coordinates)}\\
$\alpha=0^\circ$ & $1.56$ & $2.22$ & $2.68$ & $3.02$ & $3.42$ & $4.60$ & $5.35$\\
$\alpha=15^\circ$ & $2.53$ & $5.98$ & $15.5$ & $37.2$ & $130$ & $1.2\cdot10^{4}$ & $1.3\cdot10^{4}$\\
$\alpha=30^\circ$ & $9.95$ & $240$ & $1.1\cdot10^{3}$ & $1.6\cdot10^{3}$ & $6.2\cdot10^{3}$ & $4.2\cdot10^{6}$ & $1.2\cdot10^{6}$\\
$\alpha=45^\circ$$^\dagger$ & $186$ & $1.5\cdot10^{3}$ & $1.3\cdot10^{4}$ & $5.0\cdot10^{5}$ & $8.3\cdot10^{5}$ & $1.4\cdot10^{6}$ & $1.3\cdot10^{8}$\\
$\alpha=60^\circ$$^\dagger$ & $290$ & $1.1\cdot10^{4}$ & $8.5\cdot10^{4}$ & $2.8\cdot10^{5}$ & rej. & rej. & rej.\\
$\alpha=90^\circ$$^\dagger$ & $61.2$ & $3.4\cdot10^{3}$ & $8.3\cdot10^{5}$ & $2.2\cdot10^{5}$ & $3.3\cdot10^{5}$ & rej. & rej.\\
\midrule
\multicolumn{8}{l}{\PH\ (unsymmetric, hybrid space)}\\
$\alpha=0^\circ$ & $1.58$ & $2.22$ & $2.69$ & $3.00$ & $3.41$ & $3.47$ & $4.15$\\
$\alpha=15^\circ$ & $1.76$ & $2.40$ & $2.80$ & $3.10$ & $3.47$ & $4.47$ & $4.18$\\
$\alpha=30^\circ$ & $2.00$ & $4.77$ & $7.52$ & $16.0$ & $24.9$ & $74.5$ & $94.6$\\
$\alpha=45^\circ$$^\dagger$ & $2.78$ & $10.0$ & $21.2$ & $46.7$ & $91.0$ & $258$ & $546$\\
$\alpha=60^\circ$$^\dagger$ & $4.26$ & $18.9$ & $53.3$ & $193$ & $833$ & rej. & rej.\\
$\alpha=90^\circ$$^\dagger$ & $389$ & $2.9\cdot10^{3}$ & $4.8\cdot10^{3}$ & $1.7\cdot10^{3}$ & $2.2\cdot10^{3}$ & rej. & rej.\\
\midrule
GLL & $1.56$ & $2.22$ & $2.68$ & $3.00$ & $3.41$ & $3.47$ & $4.15$\\
\bottomrule
\end{tabular}
}
  \end{minipage}\hfill
  \begin{minipage}[t]{0.48\linewidth}\centering
    \captionof{table}{Ratio of the maximum element condition numbers
      $\sigma_{\max}/\sigma_{2n-3}$, \PM/\PP, $N=8$.}\label{tab:cond}
    \resizebox{\linewidth}{!}{\begin{tabular}{lrrrrrrr}
\toprule
& \multicolumn{7}{c}{$p$}\\
\cmidrule(lr){2-8}
 & 2 & 3 & 4 & 5 & 6 & 7 & 8\\
\midrule
$\alpha=0^\circ$ & 0.998 & 1.01 & 0.996 & 1 & 0.997 & 1.01 & 0.998\\
$\alpha=15^\circ$ & 1.54 & 0.937 & 1.19 & 1.61 & 3.25 & 182 & 128\\
$\alpha=30^\circ$ & 60 & 11.5 & 145 & 623 & 90.3 & 1.4e5 & 2.1e6\\
$\alpha=45^\circ$$^\dagger$ & 44.4 & 12.8 & 148 & 7.2e4 & 1.1e5 & 6.0e6 & 2.4e6\\
$\alpha=60^\circ$$^\dagger$ & 10.4 & 10.7 & 104 & 7.2e3 & 2.2e5 & 3.6e7 & 3.5e10\\
$\alpha=90^\circ$$^\dagger$ & 9.95 & 198 & 23.2 & 5.7e3 & 3.7e4 & 1.6e5 & 2.8e7\\
\bottomrule
\end{tabular}
}
  \end{minipage}
\end{table}
